\paragraph{Generic Errors}\label{sec:generic-errors}

So far we have discussed two idealized errors:
\begin{itemize}
    \item Phase-flip errors (Pauli-$Z$) that flip the phase of a qubit (\autopageref{sec:phase-flip-errors});
    \item Bit-flip errors (Pauli-$X$) that flip the state of a qubit (\autopageref{sec:bit-flip-errors}).
\end{itemize}
However, real quantum hardware is messy (i.e., noisy) and a qubit may experience electromagnetic noise, thermal fluctuations, imperfect control pulses and interactions with the environment. Therefore, the resulting \textbf{disturbance} is generally \textbf{not} a \hl{perfect phase-flip or bit-flip error}. Instead, the error can be some arbitrary transformation.

\highspace
To describe this situation, we introduce an \definition{Error Operator} $E$ that acts on the qubit exactly like a \textbf{quantum gate}:
\begin{equation}
    \ket{\psi} \rightarrow E \ket{\psi}
\end{equation}
But with the crucial difference that $E$ is \textbf{unwanted}. At first glance this \emph{seems} disastrous because there are infinitely many possible error operators. Therefore, \emph{\textbf{how could we possibly correct all of them?}}

\highspace
\begin{flushleft}
    \textcolor{Green3}{\faIcon{check-circle} \textbf{Decomposition into Pauli Operators}}
\end{flushleft}
The key insight is that \textbf{every} $2 \times 2$ complex matrix can be written as a linear combination of four special matrices:
\begin{equation*}
    I, \quad X, \quad Y, \quad Z
\end{equation*}
These matrices form a basis for the space of single-qubit operators (\autopageref{sec:pauli-gates}). So \textbf{any error operator $E$ can be written as}:
\begin{equation*}
    E = \alpha_0 I + \alpha_1 X + \alpha_2 Y + \alpha_3 Z
\end{equation*}
Where the \textbf{coefficients} $\alpha_i$ with $i = 0, 1, 2, 3$ are complex numbers that \textbf{depend on the specific physical error}.

\highspace
\begin{theorem}[Pauli Basis Decomposition]
    The set:
    \begin{equation*}
        \left\{I, X, Y, Z\right\}
    \end{equation*}
    Forms a basis for the vector space of all $2 \times 2$ complex matrices. Therefore, any single-qubit operator $E$ can be uniquely written as:
    \begin{equation}\label{eq:pauli-decomposition}
        E = \alpha_0 I + \alpha_1 X + \alpha_2 Y + \alpha_3 Z
    \end{equation}
    Where $\alpha_0, \alpha_1, \alpha_2, \alpha_3 \in \mathbb{C}$ are complex coefficients that depend on the specific operator $E$.
\end{theorem}

\newpage

\begin{flushleft}
    \textcolor{Green3}{\faIcon{question-circle} \textbf{Meaning of Each Term}}
\end{flushleft}
Each term in the decomposition has a clear physical interpretation:
\begin{itemize}
    \item \important{Identity term} ($I$): represents the case where no error occurs, and the qubit remains unchanged.
    \begin{equation*}
        I = \begin{bmatrix}
            1 & 0 \\
            0 & 1
        \end{bmatrix}
    \end{equation*}

    \item \important{Bit-flip term} ($X$): represents the probability amplitude of a bit-flip error, which flips the state of the qubit from $\ket{0}$ to $\ket{1}$ and vice versa (\autopageref{sec:bit-flip-errors}).
    \begin{equation*}
        X = \begin{bmatrix}
            0 & 1 \\
            1 & 0
        \end{bmatrix}
    \end{equation*}

    \item \important{Phase-flip term} ($Z$): represents the probability amplitude of a phase-flip error, which flips the phase of the qubit from $\ket{1}$ to $-\ket{1}$ while leaving $\ket{0}$ unchanged (\autopageref{sec:phase-flip-errors}).
    \begin{equation*}
        Z = \begin{bmatrix}
            1 & 0 \\
            0 & -1
        \end{bmatrix}
    \end{equation*}

    \item \important{Combined bit-and-phase-flip term} ($Y$): represents the probability amplitude of an error that combines both bit-flip and phase-flip effects, flipping the state of the qubit and changing its phase simultaneously, $Y = iXZ$.
    \begin{equation*}
        Y = \begin{bmatrix}
            0 & -i \\
            i & 0
        \end{bmatrix}
    \end{equation*}
\end{itemize}

\highspace
\begin{flushleft}
    \textcolor{Green3}{\faIcon{check-circle} \textbf{Why Correcting $X$ and $Z$ Errors is Sufficient}}
\end{flushleft}
Since $Y$ can be expressed as a combination of $X$ and $Z$ (i.e., $Y = iXZ$), if we can correct for $X$ and $Z$ errors, we can also correct for $Y$ errors. Therefore, \textbf{correcting for $X$ and $Z$ errors is sufficient to correct for any arbitrary error $E$}. This is a fundamental principle in quantum error correction, allowing us to design codes that can protect against a wide range of errors by focusing on correcting just these two types of errors. Depending on the source, the related theorem has many names. However, the important point is that \hl{if a quantum error-correcting code can correct all Pauli errors, then it can correct any arbitrary single-qubit error}. We will discuss quantum error-correcting codes in future sections.

\highspace
\begin{theorem}[Discretization Pauli Error, Error Discretization, Pauli Error Expansion, Digitization of Quantum Errors]
    Let \autoref{eq:pauli-decomposition}:
    \begin{equation*}
        E = \alpha_0 I + \alpha_1 X + \alpha_2 Y + \alpha_3 Z
    \end{equation*}
    Be an arbitrary error acting on a physical qubit. If a QECC (quantum error-correcting code) can correct each Pauli operator $I$, $X$, $Y$ and $Z$, individually, then it can correct the arbitrary error $E$ as well.
\end{theorem}